\chapter{Machine Learning in a nutshell} 
\label{cap:machine_learning}

Machine learning is a subfield of artificial intelligence that focuses on the development of algorithms and statistical models that enable computers to perform tasks without explicit instructions. Instead, these systems learn from data, identifying patterns and making decisions based on the information they have been trained on.
Machine learning consists of designing efficient and accurate prediction algorithms. More generally, learning techniques are data-drivem methods combining fundamental concepts in computer science with ideas from statistics, probability and optimization.
\subsubsection{Types of tasks}
The following are some standard types of tasks in machine learning:
\begin{itemize}
    \item \textbf{Classification:} Assigning labels to data points based on learned patterns (e.g., e-mail spam detection).
    \item \textbf{Regression:} Predicting continuous values based on input features (e.g., predicting house prices). In regression, the penalty for an incorrect prediction depends on the magnitude of the difference between the true and predicted values, in contrast with classification problem, where there is typically no notion of closeness between various categories. 
    \item \textbf{Object detection:} Identifying and localizing objects within images or videos (e.g., detecting tumours in lung CT scans).
    \item \textbf{Clustering:} Grouping similar data points together without predefined labels.
    \item \textbf{Anomaly detection:} Identifying unusual patterns that do not conform to expected behavior.
    \item \textbf{Ranking}: Ordering items based on their relevance or importance.
\end{itemize}

Algorithms that solve a learning task based on semantically annotated historical data are said to operate in a \textbf{supervised learning} mode. In contrast, algorithms that use data without any semantic annotation are said to operate in an \textbf{unsupervised learning} mode. In the latter case, the algorithm is expected to discover patterns in the data without any prior knowledge of the labels or categories.
In this thesis I'll mainly focus on supervised learning. 

\textbf{Label set: } We use \( Y \) to denote the set of all possible labels for a data point of a given learning problem. Note that the labels can be of two different types: \text{categorical} labels, which are discrete and finite and define classification problems, and \text{continuous} labels, which can take any value in a continuous range and define regression problems.

\section{Neural networks and deep learning} 
Neural networks are a class of machine learning algorithms inspired by the structure and function of the human brain. They consist of interconnected nodes (neurons) that process information in layers. Each connection between nodes carries a weight, which is adjusted during training to optimize performance (Figure \ref{fig:nodes}). 
Deep learning refers to the use of neural networks with many layers (deep neural networks) to model complex patterns in large datasets. This approach has led to significant advancements in areas such as image recognition, natural language processing, and game playing.

\begin{figure}[!ht]
    \centering
    \includegraphics[width=0.8\textwidth]{immagini/cap2/nodes.png}
    \caption{Schematic representation of a deep learning network. Inputs are processed through layers of interconnected nodes, where connections are weighted. The network learns by adjusting these weights to produce the desired output. }
    \label{fig:nodes}
\end{figure}

\section{YOLO: You Only Look Once}
\label{sec:yolo}
Object detection is a task that involves identifying and classifying objects present in images or videos. Initially, object detection was approached as a pipeline consisting of three main steps: proposal generation, feature extraction and region classification. However, this approach was computationally expensive and often led to suboptimal results.
The emergence of deep learning brought a significant change in object detection, with deep convolutional neural networks (CNNs) playing a crucial role in this transformation. CNNs are designed to automatically learn hierarchical features from raw pixel data, eliminating the need for manual feature engineering. This shift allowed for more efficient and accurate object detection systems.

Currently, deep learning-based object detection frameworks can be classified in two families: 
\begin{itemize}
    \item \textbf{Two-stage detectors:} These methods first generate region proposals and then classify them. Examples include R-CNNs (Region-based Convolutional Neural Networks), that first generate region proposals using a selective search algorithm and then extracts features from these regions using a CNN; the extracted features are then fed into a support vector machine for object classification. 
    \item \textbf{One-stage detectors:} These methods perform detection in a single pass, directly predicting bounding boxes and class probabilities. Examples include YOLO (You Only Look Once). The YOLO models are popular for their accuracy and compact size. It is a state-of-the-art model that could be trained on any hardware. In this thesis we use the 8th version of the software, YOLOv8  \cite{yolov8_ultralytics}, which was developed by Ultralytics and introduced on January 2023. It is used to detect objects in images, classify images and distinguish onjects from each other. 
\end{itemize}

\subsection{Architecture of YOLOv8}
The YOLOv8 architecture is composed of two major parts, namely the \textbf{backbone} and the \textbf{head}, both of which use a fully convolutional neural network, shown in Figure \ref{fig:yolo_architecture}. \cite{inbook}
\begin{figure}[!ht]
    \centering
    \includegraphics[width=0.8\textwidth]{immagini/cap2/yolo_architecture.png}
    \caption{Architecture of YOLOv8. The backbone extracts features from the input image, while the head predicts bounding boxes and class probabilities.}
    \label{fig:yolo_architecture}
\end{figure}

\begin{itemize}
    \item The \textbf{backbone} is responsible for extracting features from the input image. It consists of a modified version of the CSPDarknet53 architecture \cite{DBLP:journals/corr/abs-1911-11929}, which has 53 convolutional layers and employs a technique called cross-stage partial connections to enhance the transmission of information across the various levels of the network. The convolutional layers are organized in a sequential manner to extract relevant features from the input image.
    \item The \textbf{head} is responsible for predicting bounding boxes and class probabilities. It consists of a series of convolutional layers that take the features extracted by the backbone and apply additional operations to predict the bounding boxes and class probabilities for each object in the image. The head uses a technique called anchor boxes to handle objects of different sizes and aspect ratios.
\end{itemize}


The YOLOv8 framework can be used to perform computer vision tasks such as detection, segmentation, classification and pose estimation and comes with pre-trained models for each task. For detection, the models are pre trained on the COCO dataset, while for classification on ImageNet dataset. 
There are different versions of YOLOv8, each designed for different tasks and with different architectures. The most common versions are YOLOv8n, YOLOv8s, YOLOv8m, YOLOv8l and YOLOv8x, where the letter indicates the size of the model (n for nano, s for small, m for medium, l for large and x for extra large). The larger the model, the more parameters it has and the more computational resources it requires to train and run.

\section{Object Detection Metrics}
\label{sec:metrics}
A metric is a function that quantifies and evaluates the performance of a model. In the context of object detection, metrics are used to assess how well a model can detect and classify objects in images. Here I list the most used ones in object detection and identification problems.

\subsection*{Intersection over Union (IoU)}
When we talk about object detection, we often refer to the concept of \textbf{bounding boxes}. A bounding box is a rectangular box that is used to define the location of an object in an image, typically represented by its top-left corner coordinates (x, y), width, and height, or top-left and bottom-right coordinates. The bounding box is used to localize the object within the image and is often used in conjunction with a classification label to identify the object.
Intersection over Union (IoU) is a metric used to evaluate the accuracy of the bounding box detection. 

Two boxes are considered when evaluating IoU, the \textbf{ground truth} one, which is provided by the dataset and the \textbf{predicted} one, which is generated by the detection model. 
IoU measures the degree of overlap between the two boxes by computing the ratio between the intersection area and the union area of the ground trurh and the predicted bouding box, as shown in Figure \ref{fig:iou} and in Eq. \ref{eq:IoU}. 
Formally, if $A_{true}$ and $A_{pred}$ denote the areas of the ground truth and the predicted boxes respectively, the intersection is given by $A_{pred}\cap A_{true}$ while the union area is given by $A_{pred} \cup A_{true}$.

\begin{figure}[!ht]
    \centering
    \includegraphics[width=0.7\textwidth]{immagini/cap2/iou_2.png}
    \caption{Visual representation of the intersection over union metric. $A_{true}$ denotes the area of the ground truth bounding box, $A_{pred}$ denotes the area of the predicted bounding box. The numerator corresponds to the intersection area,$A_{pred}\cap A_{true}$, while the denominator corresponds to the union area $A_{pred}\cup A_{true}$. }
    \label{fig:iou}
\end{figure}

The IoU is calculated as follows:
\begin{equation}
    IoU = \frac{Area_{Intersection}}{Area_{Union}} = \frac{A_{pred} \cap A_{true}}{A_{pred} \cup A_{true}}
    \label{eq:IoU}
\end{equation}

The IoU value ranges from 0 to 1, where 0 indicates no overlap and 1 indicates perfect overlap. A common threshold for considering a detection as a true positive is an IoU of 0.5 or higher.

The computation of the Intersection over Union (IoU) is detailed in Algorithm \ref{alg:iou}.

\begin{algorithm}[H]
    \caption{Intersection over Union (IoU) between two bounding boxes}
    \label{alg:iou}
    \begin{algorithmic}[1]
        \State \textbf{Input:} Boxes $A = (x_A, y_A, w_A, h_A)$ and $B = (x_B, y_B, w_B, h_B)$
        
        \State \textbf{Compute bottom-right corners:}
        \Statex \hspace{\algorithmicindent} $x2_A = x_A + w_A$, \quad $y2_A = y_A + h_A$
        \Statex \hspace{\algorithmicindent} $x2_B = x_B + w_B$, \quad $y2_B = y_B + h_B$
        
        \State \textbf{Compute intersection rectangle:}
        \Statex \hspace{\algorithmicindent} $x_{\text{int}}^{1} = \max(x_A, x_B)$, \quad $y_{\text{int}}^{1} = \max(y_A, y_B)$
        \Statex \hspace{\algorithmicindent} $x_{\text{int}}^{2} = \min(x2_A, x2_B)$, \quad $y_{\text{int}}^{2} = \min(y2_A, y2_B)$

        \State \textbf{Compute intersection area:}
        \Statex \hspace{\algorithmicindent} $w_{\text{int}} = \max(0, x_{\text{int}}^{2} - x_{\text{int}}^{1})$
        \Statex \hspace{\algorithmicindent} $h_{\text{int}} = \max(0, y_{\text{int}}^{2} - y_{\text{int}}^{1})$
        \Statex \hspace{\algorithmicindent} $A_{\text{int}} = w_{\text{int}} \ast h_{\text{int}}$
        
        \State \textbf{Compute areas of each box:}
        \Statex \hspace{\algorithmicindent} $A_A = w_A \ast h_A$, \quad $A_B = w_B \ast h_B$
        
        \State \textbf{Compute union area:}
        \Statex \hspace{\algorithmicindent} $A_{\text{union}} = A_A + A_B - A_{\text{int}}$
        
        \State \textbf{Return:} $\text{IoU} = \dfrac{A_{\text{int}}}{A_{\text{union}}}$
    \end{algorithmic}
\end{algorithm}



\subsection*{Precision and Recall}
Precision and recall are fundamental metrics used to evaluate the performance of object detection models. These metrics provide valuable insights into the model's ability to identify objects of interest within images. 
Let's define some fundamental concepts before diving into precision and recall:
\begin{itemize}
    \item \textbf{True Positive (TP):} These are istances where the model currently identifies and localize objectys, and the intersection over union (IoU) score between the predicted and the ground truth bounding box is equal or greater then a specified treshold (Figure \ref{fig:TP}).
    \item \textbf{False Positive (FP):} These are cases where the model incorrectly identifies an object that does not exist in the ground truth or where the predicted box has an IoU score below tha define treshold (Figure \ref{fig:FP}).
    \item \textbf{False Negative (FN):} These rapresent instances where the model fails to detect an object that is present in the ground truth. This means that the model has not detected an object that is present in the image, resulting in a missed detection (Figure \ref{fig:FN}).
    \item \textbf{True Negative (TN):} Not applicable in object detection. It represents correcly rejecting the absence of objects, but in object detection, the goal is to detect objects rather than the absence of them.
\end{itemize}

\begin{figure}[!ht]
\centering
\begin{minipage}[t]{0.3\textwidth}
    \centering
    \includegraphics[width=\linewidth]{immagini/cap2/TP.png}
    \caption{The object is there, and the model detects it, with an IoU above treshold.}
    \label{fig:TP}
\end{minipage}
\hfill
\begin{minipage}[t]{0.3\textwidth}
    \centering
    \includegraphics[width=\linewidth]{immagini/cap2/FP.png}
    \caption{The object is there, but the predicted box has an IoU below the treshold.}
    \label{fig:FP}
\end{minipage}
\hfill
\begin{minipage}[t]{0.3\textwidth}
    \centering
    \includegraphics[width=\linewidth]{immagini/cap2/FN.png}
    \caption{The object is there, and the model doesn't detect it. The ground truth object has no prediction.}
    \label{fig:FN}
\end{minipage}
\caption{Ground truth box is the green rectangle, while the predicted box is the red rectangle.}
\label{fig:TP_FP_FN}
\end{figure}

Let's now define the core metrics: 
\begin{itemize}
    \item \textbf{Precision}: it serves to quantify the accuracy of the positive predictions made by the model. It specifically assesses how well the model distinguishes true objects from false positives. In essence, precision provides insight into the model’s ability to make positive predictions that are indeed accurate. A high precision score indicates that the model is skilled at avoiding false positives and provides reliable positive predictions.
    \begin{equation}
        Precision = \frac{TP}{TP + FP}
    \end{equation}

    \item \textbf{Recall}: also known as sensitivity or true positive rate, is another essential metric used in evaluating model performance, especially in object detection tasks. Recall measures the model’s capability to capture all relevant objects in the image. In essence, recall assesses the model’s completeness in identifying objects of interest. A high recall score indicates that the model effectively identifies most of the relevant objects in the data.
    \begin{equation}
        Recall = \frac{TP}{TP + FN}
    \end{equation}
\end{itemize}
We can define the \textbf{Average Precision (AP)} as the area under the precision-recall curve, which is obtained by plotting precision against recall for different confidence thresholds \footnote{The confidence threshold is a tunable parameter used to filter out low-confidence predictions in object detection tasks. Each predicted bounding box is assigned a confidence score, typically representing the estimated probability that the detection both contains an object and correctly identifies its class. Varying this threshold affects the trade-off between precision and recall: lower thresholds increase recall but may introduce more false positives, whereas higher thresholds improve precision at the cost of reduced recall.}. The AP is a single value that summarizes the model's performance across different levels of precision and recall. It is calculated as follows:
\begin{equation}
    AP = \int_{0}^{1} P_{\alpha}(box_{gt}, box_{pred}) dR_{\alpha}
\end{equation}
where \( P_{\alpha} \) is the precision at a given IoU threshold \( \alpha \) and \( R_{\alpha} \) is the recall at the same threshold. AP achieves a harmonious balance between false positive and false negatives, providing a comprehensive evaluation of the model's performance. 

\subsection*{Mean Average Precision (mAP)}
Mean Average Precision (mAP) extends the concept of Average Precision (AP) by providing a global summary of a model’s performance across multiple object classes and, in some cases, multiple localization thresholds. While AP evaluates the precision-recall trade-off for a single class — by varying the confidence threshold and computing the area under the resulting precision-recall curve — mAP averages these AP values to reflect the model’s overall detection capabilities.
\begin{equation}
    mAP = \frac{1}{N} \sum_{i=1}^{N} AP_i
\end{equation}
where \( N \) is the number of classes and \( AP_i \) is the Average Precision for class \( i \).

In practice, several versions of mAP are commonly reported. mAP@0.5 refers to the mean AP calculated at a fixed Intersection over Union (IoU) threshold of 0.5. mAP@[0.5:0.95], computes the average AP across ten IoU thresholds ranging from 0.5 to 0.95 (in increments of 0.05), providing a more rigorous and comprehensive evaluation that accounts for both classification and localization accuracy.

\section{Transfer learning}
Transfer learning is a machine learning technique in which knowledge gained through one task or data set is used to improve the performance of models on another related task and/or on a different data set (Figure \ref{fig:transfer_learning}). 
In other words, it uses what has been learned in one setting to improve generalization in another. The applications of transfer learning in deep learning are very interesting because models need a big amount of labeled data to learn and gain knoweledge, and a lot of times such big datasets are not avaiable.

In this thesis, transfer learning plays a central role in addressing the extreme scarcity of annotated medical images. In addition to conventional pre-training on large natural image datasets such as COCO, we investigate a novel source domain: synthetic jet images from particle physics, generated with full control over complexity and noise (Chapter \ref{cap:jets}). This approach aims to produce richer and more generalizable feature representations, ultimately improving model performance when very limited medical data is available.
\begin{figure}[!ht]
    \centering
    \includegraphics[width=0.8\textwidth]{immagini/cap2/transfer_learning.png}
    \caption{Visual representation of the difference between traditional (left) and transfer learning (right).}
    \label{fig:transfer_learning}
\end{figure}
Traditional learning processes create a new model for each new activity based on the labeled data available. This is because traditional machine learning algorithms assume that training and test data come from the same feature space, so if the data distribution changes or if the trained model is applied to a new dataset, users need to retrain a newer model from scratch, even if they want to carry out an activity similar to that of the first model. Transfer learning algorithms, on the other hand, take already trained models or networks as a starting point, then apply that model's knowledge, gained in a task or initial source data to a new, but related activity or target data.

Depending on the amount of data available and the similarity between source domain and target domain, there are different transfer learning strategies:

\begin{itemize}
    \item \textbf{Feature extraction}: Freezes most of the model, training only the head. It is useful with very little data and similar domains.
    \item \textbf{Partial fine-tuning}: Unfreezes and trains a subset of the backbone's layers (often higher layers) while freezing the rest. Suitable for moderately different domains.
    \item \textbf{Complete fine-tuning}: Unfreezes and trains all layers of the model. Suitable for very different domains or when abundant data is available.
\end{itemize}

In the context of medical imaging, transfer learning has shown promising results. Models pre-trained on large natural image datasets (e.g. ImageNet, COCO) can learn general low-level features (such as edges, textures, and shapes) that are also present in medical images. Fine-tuning such models on smaller labeled medical datasets allows for improved performance even with limited annotated data, making transfer learning a valuable tool in the medical AI pipeline.

The following chapters will explore specific cases of transfer learning applied to medical image datasets, comparing various pre-training strategies and model architectures.

\section{Curriculum Learning}

Curriculum Learning is a training strategy inspired by the way humans learn, starting from simple tasks and progressively moving to more complex ones. Introduced by Bengio et al. in 2009 \cite{inproceedings}, this approach proposes to organize the training of a model according to an ordered sequence of examples, selected based on a difficulty criterion.

Instead of presenting all data at once and in random order, Curriculum Learning involves the progressive introduction of data subsets, starting from the easiest examples and gradually increasing the complexity. This allows the model to first acquire robust and general patterns, and then refine its predictive ability by tackling more challenging cases.

\begin{figure}[!ht]
\centering
\includegraphics[width=0.8\textwidth]{immagini/cap2/curr.png}
\caption{Illustration of the curriculum learning concept \cite{9392296}.}
\label{fig:curriculum}
\end{figure}

An effective curriculum typically requires three main components:
\begin{itemize}
\item \textbf{Difficulty criterion}: the mechanism through which a difficulty level is assigned to each example, based on known characteristics (e.g., geometric properties, data quality, object dimensions) or on model-driven indicators such as early-loss values or prediction uncertainty.
\item \textbf{Pacing}: the schedule that controls how quickly more difficult examples are introduced. This can be predefined (e.g., adding a fixed number of examples per epoch) or adaptive (e.g., based on validation performance).
\item \textbf{Presentation order}: the sequence in which examples or subsets are fed to the model, which can be fixed or dynamically updated during training.
\end{itemize}

This approach has been shown to be useful in scenarios with high variability, noisy or unbalanced data, or when faster convergence is desirable. It can also help avoid the premature learning of spurious correlations or statistical noise, thus improving generalization.

\subsection*{Application in this thesis}
In this work, Curriculum Learning is enabled by the \textit{ad hoc} construction of synthetic datasets from particle physics, specifically jet images from high-energy collisions, organized into multiple difficulty levels with controlled complexity and noise. This structure allows us to systematically implement a curriculum in the \textit{source domain}, gradually increasing the challenge before transferring the learned features to the medical imaging domain. By doing so, the model is encouraged to develop richer and more generalizable representations, ultimately improving performance when fine-tuned on extremely small annotated CT datasets.

\subsection*{Practical considerations}
The effectiveness of Curriculum Learning depends heavily on:
\begin{itemize}
\item A well-defined difficulty criterion;
\item An appropriate pacing schedule;
\item The alignment of the curriculum with the target task.
\end{itemize}

If these elements are poorly designed, Curriculum Learning may fail to provide benefits and can even harm performance compared to standard training. However, when implemented carefully, it can be a powerful tool to guide model optimization, especially in cases with uneven data distributions, low signal-to-noise ratio, or tasks that benefit from gradual learning.

\newpage

\section{YOLO dataset}
Before discussing the results, it is necessary to provide some background on Ultralytics, the company that developed the YOLO model.

In this thesis, we focus on supervised training, in which the model is provided with images and corresponding annotations during the training phase. These annotations include the bounding box coordinates (top-left coordinate, height, and width) as well as the object class, since our task is object detection.

Ultralytics requires the dataset to be organized in a specific directory structure. The dataset must be split into training, validation, and test sets, as in any standard workflow. Images are stored in their respective subfolders under \texttt{images}, while annotations are placed in corresponding subfolders under \texttt{labels}. Each annotation is saved in a \texttt{.txt} file with the same name as its corresponding image.

A schematic representation of the dataset organization is shown in Figure~\ref{fig:yolo_organization}.


\begin{figure}[!ht]
    \centering
\begin{forest}
for tree={
    font=\ttfamily,
    grow'=0,
    child anchor=west,
    parent anchor=south,
    anchor=west,
    calign=first,
    edge path={
      \noexpand\path [draw, \forestoption{edge}]
      (!u.south west) +(5pt,0) |- (.child anchor)\forestoption{edge label};
    },
    before typesetting nodes={
      if n=1
        {insert before={[,phantom]}}
        {}
    },
}
[dataset
  [images
    [train]
    [val]
    [test]
  ]
  [labels
    [train]
    [val]
    [test]
  ]
]
\end{forest}
    \caption{Schematic organization of a YOLO dataset.}
    \label{fig:yolo_organization}
\end{figure}

Inside the dataset folder, a configuration file named \texttt{dataset.yaml} must be included. This file is provided to the model during training. It specifies the relative paths to the training, validation, and test image directories (the paths to the labels are inferred automatically, since they follow the same folder structure as the images, and the annotation files share the same names as their corresponding images). The file must also define the number of classes and their names.

\subsection*{Annotations}
YOLO also requires annotations to be provided in a specific format:

\begin{equation*}
    \text{class} \qquad \dfrac{x_{\text{center}}}{\text{image width}} \qquad \dfrac{y_{\text{center}}}{\text{image height}} \qquad \dfrac{\text{width}}{\text{image width}} \qquad \dfrac{\text{height}}{\text{image height}}
\end{equation*}

Each annotation file contains one line for each object present in the corresponding image. The first value is the class index, followed by the normalized coordinates of the bounding box: the x- and y-coordinates of the center of the bounding box, and its width and height, all normalized by the dimensions of the image.

\begin{figure}[!ht]
    \centering
    \includegraphics[width=0.55\textwidth]{immagini/cap6/annotations.png}
    \caption{Example of YOLO annotation.}
    \label{fig:annotations_cat}
\end{figure}


\section{Training Settings}
Training a deep learning model involves feeding it data and adjusting its parameters so that it can make accurate predictions. 
\subsection{Hyperparameters}
The training process is controlled by a set of hyperparameters, which are parameters that are not learned during training but are set before the training begins. These hyperparameters include:
\begin{itemize}
    \item \textbf{Learning Rate}: This is a crucial hyperparameter that determines how much the model's weights are updated during training. A higher learning rate can lead to faster convergence but may also cause the model to overshoot the optimal solution.
    \item \textbf{Batch Size}: This refers to the number of training examples used in one iteration of training. A larger batch size can lead to more stable gradients but requires more memory.
    \item \textbf{Number of Epochs}: An epoch is one complete pass through the entire training dataset. The number of epochs determines how many times the model will see the entire dataset during training.
    \item \textbf{Optimizer}: The optimizer is an algorithm used to update the model's weights based on the computed gradients. Common optimizers include Adam, SGD.
    \item \textbf{Data Augmentation}: This technique involves applying random transformations to the training data to increase its diversity and improve the model's generalization ability. Common augmentations include rotation, scaling, flipping, and color adjustments.
    \item \textbf{Early Stopping Criteria}: This is a condition that stops training when the model's performance on the validation set does not improve for a specified number of epochs. It helps prevent overfitting and saves computational resources.
    \item \textbf{Weight Decay}: technique used to prevent overfitting by adding a penalty term to the loss function based on the magnitude of the model's weights. It encourages the model to learn simpler patterns.
    \item \textbf{Dropout}: A regularization technique that randomly sets a fraction of the model's neurons to zero during training. This helps prevent overfitting by reducing the model's reliance on specific neurons.
\end{itemize}

\subsection{Loss Function}
The loss function is a critical component of the training process, as it quantifies how well the model's learning and the predictions match the ground truth. In the case of object detection using YOLOv8, the loss function typically consists of three components:

\subsubsection{1. Box Loss}
The box loss measures the difference between the predicted bounding box coordinates and the ground truth coordinates. It is calculated as the Complete Intersection over Union (\textbf{CIoU}) loss, which is an extension of the Intersection over Union metric discussed in Chapter 2, and takes into account not only the overlap between the predicted and ground truth boxes but also their aspect ratio and distance between their centers. The CIoU loss is defined as:
\begin{equation}
    L_{CIoU} = 1 - IoU - \dfrac{d^2}{c^2} - \alpha v
\end{equation}
where IoU is defined in Eq.\ref{eq:IoU}, where $v$ is a function of the boxes widths and heights and $\alpha$ is a trade-off parameter function of IoU. 

\begin{figure}[!ht]
    \centering
    \includegraphics[width=0.5\textwidth]{immagini/cap6/CIoU.png}
    \caption{Representation of a ground truth box in green and its prediction in red. "d" is the euclidean distance between the centers, while "c" is the diagonal of the smallest box including the green and the red one.}
    \label{fig:CIoU}
\end{figure}

\subsubsection{2. Classification Loss}

\subsubsection{3. Dual Focal Loss}
The classification loss measures the difference between the predicted class probabilities and the ground truth class labels. YOLOv8 uses a dual focal loss, which is a combination of binary cross-entropy loss and focal loss. The binary cross-entropy loss is used for multi-label classification tasks, while the focal loss helps to address class imbalance by down-weighting easy-to-classify examples and focusing more on hard-to-classify examples.
The dual focal loss is defined as:
\begin{equation}
    L_{dual} = \dfrac{1}{N} \sum_{i=1}^{N} \left( -\alpha (1 - p_i)^{\gamma} \log(p_i) - (1 - \alpha) p_i^{\gamma} \log(1 - p_i) \right)
\end{equation}
where $N$ is the number of classes, $p_i$ is the predicted probability for class $i$, $\alpha$ is a balancing factor, and $\gamma$ is a focusing parameter that controls the down-weighting of easy examples.
\subsubsection{Total Loss}
The total loss for YOLOv8 is a weighted sum of the box loss, classification loss, and dual focal loss:
\begin{equation}
    L_{total} = \lambda_{box} L_{CIoU} + \lambda_{cls} L_{cls} + \lambda_{dual} L_{dual}
\end{equation}
where $\lambda_{box}$, $\lambda_{cls}$, and $\lambda_{dual}$ are hyperparameters that control the relative importance of each loss component and can be set in the configuration scrpit. Default values are: $\lambda_{box} = 7.5$, $\lambda_{cls} = 0.5$, and $\lambda_{dual} = 1.5$.

\subsection{Augmentation}
Augmentation techniques are essential for improving the robustness and performance of YOLO models by introducing variability into the training data, helping the model generalize better to unseen data. Here's a list of the augmentations that have been used in this work and can be set in the configuration file with the others hyperparameters:
\begin{itemize}
    \item \texttt{hsv}: modifies the hue, saturation, and value of the image, allowing the model to learn to recognize objects under different lighting conditions and color variations.
    \item \texttt{degrees}: applies random rotations to the image, helping the model become invariant to object orientation. 
    \item \texttt{translate}: shifts the image horizontally and/or vertically, allowing the model to learn to recognize objects that may not be centered in the frame.
    \item \texttt{scale}: resizes the image, helping the model learn to recognize objects at different scales.
    \item \texttt{shear}: applies a shear transformation to the image, which can help the model learn to recognize objects that may be distorted or skewed.
    \item \texttt{flipud}: flips the image vertically, allowing the model to learn to recognize objects that may appear upside down.
    \item \texttt{fliplr}: flips the image horizontally, allowing the model to learn to recognize objects that may appear mirrored.
    \item \texttt{mosaic}: combines multiple images into a single image, creating a more diverse training set and helping the model learn to recognize objects in different contexts.
    \item \texttt{mixup}: combines two images and their corresponding annotations, creating a new image that contains features from both. This technique helps the model learn to recognize objects that may be partially occluded or overlapping.
\end{itemize}
